
Following Section 2.3 of  {\it Discrete Mathematics, an Open Introduction}.  Let $\{a_n\}$ be a sequence of real numbers.  The {\bf $k^{th}$ forward difference} of this sequence are defined recursively:\\
The {\bf first forward difference} $\Delta a_n$ is an new sequence defined by:$\Delta a_n=a_{n+1}-a_{n}$.\\
The {\bf $k$th forward difference} $\Delta^{k} a_n$ is $\Delta^{k} a_n=\Delta^{k-1} a_{n+1}- \Delta^{k-1}a_{n}$.  So if the sequence for $a_n$ is $1,3,5,7,9, \ldots$ then  $\Delta a_n$ is $3-1,5-3,7-5,9-7, \ldots= 2,2,2,2,2, \ldots$.

\begin{exercise}
Find $\Delta$ and $\Delta^2$ for the following sequences :\\
a. $2, 4, 6, 8, 10, 12, 14, \ldots$\\
b.  $1, 4, 9, 16, 25, 36, 49, \ldots$
\end{exercise}

\begin{exercise}
Make up sequences that have:\\
a. 3, 3, 3, 3, … as its second differences.\\
b. 1, 2, 3, 4, 5, … as its third differences.\\
c. 1, 2, 4, 8, 16, … as its fifth differences.\\
\end{exercise}

%\begin{exercise}
%Using the relationships in Proof 14.1 and  $a_{n-3}=a_n - 3\Delta a_n + 3\Delta^2 a_n-\Delta^3 a_n$, rewrite $a_n = 3a_{n-1} - 3a_{n-2} + a_{n-3}$ in terms of $a_n$, $\Delta a_n$, $\Delta^2 a_n$, and $\Delta^3 a_n$.
%\end{exercise}


\begin{exercise}
Find the closed form solution to the sequence using the using the polynomial fitting method for the sequence $0, 3, 8, 15, 24, 35, 48, 63, 80,\ldots$ which has the recurrence $a_n = a_{n-1} + 2n+1$ with $a_0=0$.
\end{exercise}

\begin{exercise}
Find the closed form solution to the sequence using the using the polynomial fitting method for the sequence $1, 4, 10, 20, 35, 56, 84,\ldots$ which has the recurrence $a_n = 4a_{n-1} - 6a_{n-2} + 4a_{n-3} - a_{n-4}$ with $a_1=1, a_2=4, a_3=10$ and $a_4=20$.  This also happens to the be the sequence given by summing up the triangular numbers.  That is, $a_n=\sum_{i=1}^n T_n$.
\end{exercise}

\begin{exercise}
Calculate $\Delta \binom{n}{4}$.  Use this to rewrite $\Delta \binom{n}{k}$ as a binomial coefficient.
\end{exercise}

%\begin{proofp}
%Show $a_{n-1}=a_n - \Delta a_n$ and $a_{n-2}=a_n - 2\Delta a_n + \Delta^2 a_n$
%\end{proofp}

\begin{proofp}
Prove $\Delta (\Delta^i a_n)=\Delta^{i+1} a_n$.
\end{proofp}

\begin{proofp}
Let $b_n=n^3$.  Prove $\Delta^2 b_n=6n+6$.
\end{proofp}

\begin{proofp}$\bigstar$
Let $b_n=n^k$.  Prove $\Delta b_n=\sum_{i=0}^{k-1}\binom{k}{i}n^{i}$
\end{proofp}

\begin{proofp}$\bigstar$
Let $a_n=\alpha^n$ for $n \ge 0$ and $\alpha$ a real number.  Prove $\Delta^k a_n=(\frac{\alpha-1}{\alpha})^k a_n$ for all $k \ge 0$.
\end{proofp}

\begin{proofp}$\bigstar$
Let $a_n$ and $b_n$ be sequences.  Prove $\Delta (a_n\cdot b_n)=a_{n+1}\Delta b_n + b_{n}\Delta a_n$.
\end{proofp}

%\begin{proofp}$\bigstar$
%Given two sequences $a_n$ and $b_n$.  Prove $\Delta(a_n\cdot b_n)= a_n\cdot \Delta(b_n)+ %bigtriangledown(a_n)\cdot b_n$.
%\end{proofp}

%\begin{proofp}$\bigstar$
%Prove $a_{n-k}=\sum_{i=0}^{k}(-1)^k\binom{k}{i}\Delta^i a_{n}$. (Hint: do induction on k, and this is the generalization of the first proof above.
%\end{proofp}

%\begin{proofp}$\bigstar$
%Using the preceding proof. Prove that you can rewrite any linear homogeneous recurrence relation as a linear combination of $\Delta^i a_n$ where $\Delta^0 a_n=a_n$. This expression is call the difference equation version of the recurrence relation.
%\end{proofp}



